%% TeXora template: Tagged Article. Not slides: an article that LaTeX itself tags for screen readers
%% (PDF/UA), with nothing from TeXora in it, so it compiles the same in any editor: TeXora, Overleaf,
%% TeXstudio, TeXShop, VS Code with LaTeX Workshop, or the command line.
%%
%% What it needs:
%%   - LaTeX from 2023 or later (TeX Live 2023+, MiKTeX, Overleaf). From LaTeX 2025-06 the first
%%     \DocumentMetadata line is used (tagging=on, PDF/UA-2); before, the second (the test phase).
%%   - The picture media/test-card.png, which comes with this template.
%%   - LuaLaTeX gives the best tags (also MathML for the maths); pdfLaTeX tags too. In Overleaf:
%%     Menu > Compiler > LuaLaTeX. With LuaLaTeX, LaTeX may also write <name>-luamml-mathml.html.
%%
%% To test it in another editor: open this file there, compile, then check the PDF with
%% TeXora (Tools > Check accessibility), veraPDF or PAC, or the LaTeX project's show-pdf-tags.
%% Change the colours on the \definecolor lines below.
% !TEX program = lualatex
\IfFormatAtLeastTF{2025-06-01}
  {\DocumentMetadata{lang=en-US,pdfstandard=ua-2,tagging=on}}
  {\DocumentMetadata{lang=en-US,pdfstandard=ua-1,testphase={phase-III}}}
\documentclass[11pt]{article}
\usepackage[T1]{fontenc}
\usepackage{lmodern}
\usepackage[english]{babel}
\usepackage[margin=2.5cm]{geometry}
\usepackage{graphicx}
\usepackage{booktabs}
\usepackage{amsmath}
\usepackage{xcolor}
\usepackage{hyperref}
\definecolor{linkblue}{RGB}{0,110,200}
\definecolor{rulegrey}{RGB}{90,90,90}
\hypersetup{colorlinks,linkcolor=linkblue,urlcolor=linkblue,citecolor=linkblue,pdftitle={A Tagged Article},pdfauthor={Your Name}}

\title{A Tagged Article}
\author{Your Name\\ Your Institution}
\date{\today}

\begin{document}
\maketitle

\begin{abstract}
  This article is a test and a guide in one. Each section uses one kind of content, the way that
  keeps it readable with a screen reader: headings, lists, a described picture, a table with a
  header row, maths, links, a footnote and references. Compile it in any LaTeX editor and check the
  tags of the PDF.
\end{abstract}

\section{Headings and paragraphs}
Sections become headings (H1, H2, \dots) in the tags, so a screen reader can list them and jump
between them. Write \verb|\section|, \verb|\subsection| and \verb|\paragraph| rather than bold
text in a larger size: those carry no structure.

\subsection{Lists}
Lists are tagged as lists, item by item:
\begin{itemize}
  \item a bulleted list, with \verb|itemize|;
  \item a numbered one, with \verb|enumerate|:
    \begin{enumerate}
      \item first,
      \item second;
    \end{enumerate}
  \item and terms with their meaning, with \verb|description|.
\end{itemize}

\section{A described picture}
Figure~\ref{fig:card} has alt text, given with the \texttt{alt} key of \texttt{\textbackslash includegraphics}.
Its caption is read as well, so the alt text says what the picture shows and the caption what it
means.

\begin{figure}[htbp]
  \centering
  \includegraphics[width=0.45\textwidth,alt={A television test card: eight vertical colour bars, a rainbow strip and a frame counter}]{media/test-card.png}
  \caption{A test card, used to check that colours come through.}
  \label{fig:card}
\end{figure}

A picture that is only decoration takes \texttt{artifact} instead: \texttt{\textbackslash includegraphics[artifact]\{...\}}.

\section{A table with a header row}
Table~\ref{tab:engines} has one header row. With tagging on, LaTeX marks the first row as headers
when it is told so; the rule under it shows the same to the eye.
\IfFormatAtLeastTF{2025-06-01}{\tagpdfsetup{table/header-rows={1}}}{}

\begin{table}[htbp]
  \centering
  \caption{How each engine tags a document.}
  \label{tab:engines}
  \begin{tabular}{@{}lll@{}}
    \toprule
    Engine & Tags & Maths \\
    \midrule
    LuaLaTeX & complete & as MathML too \\
    pdfLaTeX & complete & as text \\
    XeLaTeX & incomplete & as text \\
    \bottomrule
  \end{tabular}
\end{table}

\section{Maths}
Inline maths such as $a^2 + b^2 = c^2$ and displayed equations are tagged as formulas:
\begin{equation}
  \int_0^1 x^n \, dx = \frac{1}{n+1}, \qquad n \ge 0.
\end{equation}
With LuaLaTeX the PDF also holds each formula as MathML, which screen readers can speak.

\section{Links, notes and references}
Link text should say where it goes: the
\href{https://www.w3.org/WAI/WCAG22/quickref/}{WCAG 2.2 quick reference}, not ``click
here''.\footnote{Footnotes are tagged as notes and linked to where they are called.}
References are tagged too~\cite{latex-tagging}.

\begin{thebibliography}{9}
  \bibitem{latex-tagging} The \LaTeX{} Project, \emph{Tagging status of LaTeX packages and
  classes}, \url{https://latex3.github.io/tagging-project/tagging-status/}.
\end{thebibliography}

\end{document}
